\section{Проєктування системи керування ключами }

\begin{table}[H]
\centering
\renewcommand{\arraystretch}{1.5}
\begin{tabular}{|p{2.1cm}|p{2.1cm}|p{2.1cm}|p{2.1cm}|p{2.1cm}|p{2.1cm}|}
\hline
\textbf{Секрет} & \textbf{Призна\-чення} & \textbf{Збері\-гання} & \textbf{Хто має доступ} & \textbf{Експорт} & \textbf{Ротація} \\
\hline
Ключ шифрування БД & Шифрува\-ння та розшифрування конфіденційних даних користувачів & Захищене сховище ключів (KMS) або HSM & Backend-сервіс через спеціальний механізм доступу & Ні, бажано не експортувати у відкритому вигляді & Періодич\-на заміна зі збереженням можливості розшифрування старих даних \\
\hline
Ключ цифрового підпису & Створення цифрових підписів документів та підтвердження їх автентичності & HSM або захищене сховище ключів & Backend-сервіс, який виконує підписання документів & Ні, приватний ключ не повинен експортуватися & Періодич\-на заміна ключа; старі ключі зберігаються для перевірки раніше створених підписів \\
\hline
TLS private key & & & & & \\
\hline
API credential & & & & & \\
\hline
\end{tabular}
\caption{Характеристики різних типів секретної інформації}
\label{tab:secret_types}
\end{table}
